% Section: Conclusion
\section{Conclusion}
\label{sec:conclusion}

% TODO: Expand this section.
%
% Key points to summarize:
% - The recursive drift problem in AI-augmented engineering
% - Reflector's four core mechanisms
% - The importance of governance as a first-class architectural concern
% - The path forward for reflective development systems

The increasing deployment of autonomous AI agents in software engineering workflows
demands a corresponding maturation of governance frameworks. Without explicit
mechanisms for bounding autonomy, validating state, and maintaining human alignment,
recursive AI development systems are vulnerable to drift patterns that compound
silently across development cycles.

This paper introduced the Reflector framework, which addresses this challenge through
four complementary mechanisms: scoped autonomous agents, human-in-the-loop governance
checkpoints, recursive auditing systems, and milestone synchronization protocols.
Together, these mechanisms provide a structured governance substrate on which
AI-augmented development systems can operate with greater confidence and reduced risk.

The central insight of Reflector is that autonomy and governance are not in opposition.
Well-designed governance structures \textit{enable} greater autonomous operation by
establishing the trust and verification mechanisms necessary for humans to delegate
meaningfully. The goal is not to constrain agents, but to create the conditions under
which agent autonomy can be extended safely and reversibly.

We believe that reflective development systems—systems designed from the outset to
monitor, evaluate, and govern their own recursive operation—represent a necessary
evolution in the discipline of AI-augmented software engineering.
